\begin{frame}{Conclusiones y referencias}
\scriptsize
\begin{columns}[T,totalwidth=\textwidth]
\column{0.55\textwidth}
\begin{enumerate}
  \item Los ratios transforman partidas contables en indicadores comparables.
  \item La liquidez estima pago de corto plazo, revisando calidad de activos.
  \item La gestión muestra eficiencia de cobranza e inventarios.
  \item La solvencia revela el peso del financiamiento de terceros.
  \item La rentabilidad relaciona utilidad con ventas, activos y patrimonio.
  \item Ningún ratio debe interpretarse de manera aislada.
  \item La automatización mejora cálculo y seguimiento, pero no reemplaza el análisis profesional.
\end{enumerate}

\column{0.42\textwidth}
\casehead{Referencias APA 7}

Amat, O. (2008). \textit{Análisis de estados financieros}. Gestión 2000.

Horngren, C. T., Sundem, G. L., Elliott, J. A., \& Philbrick, D. R. (2010). \textit{Introducción a la contabilidad financiera}. Pearson.

Warren, C. S., Reeve, J. M., \& Duchac, J. E. (2016). \textit{Contabilidad financiera}. Cengage Learning.

IFRS Foundation. (2018). \textit{Conceptual Framework for Financial Reporting}.
\end{columns}
\end{frame}
